\subsection{The pair attached to a specified representation category}
The coefficient field alone does not select a higher-degree replacement for the sign-Tate pair. A chosen semisimple representation category does give one explicitly.

Let $G$ be a finite group and $I$ its set of irreducible complex representations, with representatives $V_i$ and trivial representation $V_1$. Write
\[
 V_i\otimes V_j\cong\bigoplus_{k\in I} V_k^{\oplus N_{ij}^k},
 \qquad N_{ij}^k\in\N.
\]
Let $Q_G$ be the finite subsets of $I\times\Z$, with addition given by union, zero the empty set, and multiplication
\[
 S*T=\{(k,a+b): (i,a)\in S,\ (j,b)\in T,\ N_{ij}^k>0\}.
\]
Its unit is $\{(1,0)\}$. For $S\in Q_G$ put $F_G(S)=\Z^S$, with extension by zero along inclusions. The convolution of two coefficient vectors is
\[
 (c*d)_{k,n}=\sum_{a+b=n}\sum_{i,j\in I}N_{ij}^k c_{i,a}d_{j,b}.
\]
Only indices in their finite supports contribute.

\begin{proposition}[Fusion support and arithmetic degree]\label{prop:fusion-pair}
The operations above make $Q_G$ an idempotent semiring and $F_G$ an abelian-group-valued lax monoidal functor of the kind in \citet[Theorem 8.5]{CFMM}. Its inverse semiring consists of pairs $(S,c)$, $c\in\Z^S$, with the displayed convolution and support multiplication.

Let $\mathcal R_G^+$ be the semiring of finite direct sums of the symbols $V_iq^a$, using tensor decomposition for multiplication. The effective parity semiring
\[
 \mathcal E_G=\mathcal R_G^+[\eps]/(\eps^2-1)
\]
has a surjective homomorphism to this inverse semiring: at every index $(i,a)$ it records whether the sum of the two multiplicities is positive, and records their difference. An occupied coefficient may thus be zero.

For arithmetic realizations, choose an $\ell$-adic realization of the finite-group representations and an embedding of their algebraic character values into $\C$; the Tate character below is interpreted $\ell$-adically. Suppose in addition that a continuous surjection
$\operatorname{Gal}(\overline\Q/\Q)\twoheadrightarrow G$ is specified, $q$ is realized by the Tate character with unramified Frobenius eigenvalue $p$, and $\eps$ by a specified quadratic character. Then the irreducible effective constituents above an occupied $(i,a)$ have rank $\dim V_i$. At unramified primes their local factors have that degree. All such constituents have rank one if and only if every occupied $i$ is one-dimensional. This statement concerns the effective parity sectors, including occupied zero coefficients, rather than only their virtual signed difference.
\end{proposition}
\begin{proof}
Complete reducibility follows by averaging a projection over the finite group. Associativity and commutativity of tensor product give the corresponding equalities between sums of products of the nonnegative integers $N_{ij}^k$. Such a sum is positive exactly when at least one summand is positive. Taking supports of these identities proves associativity and commutativity of $*$, and distributivity over union. The trivial representation gives the claimed unit.

The convolution formula is bilinear and is compatible with extension by zero. Its associativity follows from the same tensor multiplicity identities, now as integer equalities; the coefficient of $V_1q^0$ supplies the unit constraint. These are precisely the lax monoidal identities. In the pair description, addition takes the union of supports and adds the extended vectors; multiplication is $(S,c)(T,d)=(S*T,c*d)$. The additive inverse is $(S,-c)$, so the additive monoid is inverse.

For multiplicities $p_{i,a},n_{i,a}\ge0$, put $u=p+n$ and $c=p-n$. Nonnegative convolution satisfies $\operatorname{supp}(u*v)=\operatorname{supp}(u)*\operatorname{supp}(v)$ because each $N_{ij}^k$ is nonnegative. The two sign evaluations give the usual sum and convolution identities for $u$ and $c$. Hence the stated map is a semiring homomorphism. To lift $(S,c)$, use $(p,n)=(c,0)$ where $c>0$, $(0,-c)$ where $c<0$, and $(1,1)$ at each occupied zero; outside $S$ use $(0,0)$. This proves surjectivity.

Inflation along a surjection preserves irreducibility: a subspace invariant under the Galois action is invariant under all of $G$. Tensoring with a character is an autoequivalence and preserves both irreducibility and dimension. Tate twists have the same property. An unramified local factor of $V_iq^a$ is
\[
 \det(1-p^a\rho_i(\operatorname{Frob}_p)T)^{-1};
\]
its denominator has degree $\dim V_i$ because Frobenius is invertible. A direct sum is a sum of rank-one constituents precisely when all its irreducible summands are one-dimensional. This gives the last assertion. At a ramified prime the degree is instead the dimension of the inertia-invariant subspace, so the unramified qualification is necessary.
\end{proof}

These are ordinary Artin local factors after the realization has been given; direct sums and their $L$-function factorizations are described by \citet[Remark 1.4, p.\ 203]{Sutherland}. Splitting a single local polynomial over $\C$ is not a decomposition into global rank-one representations. The latter must be compatible with every Galois element.

Two elementary examples show why rationality of coefficients cannot determine this degree. The quartic Dirichlet character modulo $5$ defined by $\chi(2)=i$ has values in $\Q(i)$ and rank one. In contrast the permutation representation of $S_3$ on $\Q^3$ has an invariant complement
\[
 V=\{(x_1,x_2,x_3):x_1+x_2+x_3=0\}
\]
of dimension two, with rational character $(2,0,-1)$ on the identity, transpositions, and three-cycles. It is irreducible over $\C$: a one-dimensional character is trivial on three-cycles, whereas the three-cycle on $V$ has the two primitive cube roots as eigenvalues. The splitting field of $X^3-2$ realizes this $S_3$ group, since the polynomial is Eisenstein at $2$ and its discriminant $-108$ is not a square. Thus rational characters already occur at rank two. Tensoring this representation with the quartic character gives a rank-two example with the same nonrational character field $\Q(i)$ as the rank-one example. Indeed, the splitting field of $X^3-2$ has unique quadratic subfield $\Q(\sqrt{-3})$, whereas the cyclic quartic field $\Q(\zeta_5)$ has quadratic subfield $\Q(\sqrt5)$. Their intersection is therefore $\Q$: a common quotient of $S_3$ and $C_4$ would otherwise be quadratic. The compositum has group $S_3\times C_4$, and an element trivial on the first factor with $\chi=i$ has trace $2i$. Thus the character field is exactly $\Q(i)$.

For the finite-image representations just used, unramified Frobenius eigenvalues are roots of unity, so the Tate-twisted local eigenvalues have absolute value $p^a$. This local purity follows directly from finite order. A statement that all nontrivial global zeros lie on the corresponding critical line is a separate Artin or Hecke Riemann hypothesis. The specified Dirichlet input addresses rank-one characters over $\Q$; it does not by itself establish that higher-rank assertion. For more general arithmetic categories one must provide their representations, tensor rules, and weight theorem first. An algebraic coefficient over $\Q(D)$ supplies none of those choices uniquely.
